\documentclass[12pt]{article}
\usepackage[margin=0.72in]{geometry}
\usepackage{fontspec}
\setmainfont{Latin Modern Roman}
\usepackage{microtype}
\usepackage{fancyhdr}
\usepackage{titlesec}
\usepackage{enumitem}
\usepackage{needspace}
\usepackage{xcolor}
\usepackage[hidelinks]{hyperref}
\setlength{\parindent}{0pt}
\setlength{\parskip}{0.34em}
\setlength{\headheight}{15pt}
\titleformat{\section}{\large\bfseries\scshape}{}{0pt}{}
\titleformat{\subsection}{\normalsize\bfseries}{}{0pt}{}
\pagestyle{fancy}
\fancyhf{}
\fancyhead[L]{Whately Logic: Week 2 Test}
\fancyhead[R]{The Two-Faced Word}
\fancyfoot[C]{\thepage}
\renewcommand{\headrulewidth}{0.4pt}

\newcommand{\focusbox}[1]{\par\vspace{0.35em}\noindent\begingroup\setlength{\fboxsep}{6pt}\colorbox{gray!12}{\begin{minipage}{\dimexpr\linewidth-2\fboxsep\relax}#1\end{minipage}}\endgroup\par\vspace{0.45em}}
\newcommand{\studentline}{\par\vspace{0.60em}\hrule\vspace{0.64em}}
\newcommand{\shortline}{\par\vspace{0.38em}\hrule\vspace{0.48em}}
\newcommand{\twomeaningsblock}{%
\textbf{Two-faced word:}\shortline
\textbf{Meaning in the first use:}\shortline
\textbf{Meaning in the second use:}\shortline
\textbf{Is this ambiguity or equivocation? Explain.}\studentline
\textbf{Rewrite the argument so the two-faced word has one clear meaning:}\studentline\studentline
\textbf{After your rewrite, does the conclusion follow? Why or why not?}\studentline
}

\begin{document}
\begin{center}
{\LARGE\bfseries Week 2 Logic Test}\par\vspace{0.16em}
{\large\scshape The Case of the Two-Faced Word}\par\vspace{0.25em}
{A Whately Puzzle on Terms, Ambiguity, and Equivocation}\par\vspace{0.45em}
\rule{0.82\textwidth}{0.4pt}
\end{center}

\section*{Name: \rule{2.35in}{0.4pt} \hfill Date: \rule{1.25in}{0.4pt}}

\focusbox{\textbf{Whately's warning:} Before asking whether a conclusion follows, ask whether the important words keep the same meaning. An argument may look strong only because one term has been allowed to carry two ideas at once. Fix the term; then test the reasoning.}

\section*{The Puzzle Story}
At the Academy of Straight Thinking, an old cabinet stands in the corner of the library. Inside it rests the \textbf{Golden Key}, the prize given to the student who can unlock a confused argument.

On the cabinet door is written:

\focusbox{\centering ``The Golden Key belongs to the student who can prove that the rule is just.''}

Beneath the words are three folded notices from the old Head Tutor. Each notice contains a trap: one word turns its face in the middle of the argument.

\textbf{The Notice of the Fair Judge says:}
\begin{quote}
``The trial was supposed to be fair. But the sky was not fair that day, for it rained. Therefore the trial was not fair.''
\end{quote}

\textbf{The Notice of the Free Hall says:}
\begin{quote}
``All students should have free speech in the Hall. Therefore no one should ever have to pay for a speech in the Hall.''
\end{quote}

\textbf{The Notice of the True Heir says:}
\begin{quote}
``Only the true heir may open the cabinet. Nora spoke truly about the cabinet. Therefore Nora is the true heir.''
\end{quote}

Three students hurry forward. Mira points to the Fair Judge notice. Elias points to the Free Hall notice. Nora points to the True Heir notice. Each claims that the notice proves the answer.

The Head Tutor has left one final sentence:

\focusbox{``Do not ask first who sounds clever. Ask first whether the word stayed faithful to one meaning.''}

\clearpage

\section*{Part I: Find and Fix the Two-Faced Word}
For each notice, identify the term that changes meaning. Then rewrite the argument so that the term has one clear meaning before you test the conclusion.

\Needspace{0.30\textheight}
\subsection*{1. The Notice of the Fair Judge}
\textbf{Argument:} ``The trial was supposed to be fair. But the sky was not fair that day, for it rained. Therefore the trial was not fair.''
\twomeaningsblock

\Needspace{0.30\textheight}
\subsection*{2. The Notice of the Free Hall}
\textbf{Argument:} ``All students should have free speech in the Hall. Therefore no one should ever have to pay for a speech in the Hall.''
\twomeaningsblock

\Needspace{0.30\textheight}
\subsection*{3. The Notice of the True Heir}
\textbf{Argument:} ``Only the true heir may open the cabinet. Nora spoke truly about the cabinet. Therefore Nora is the true heir.''
\twomeaningsblock

\Needspace{0.42\textheight}
\section*{Part II: Strengthen One Argument}
Choose one notice from Part I. Add one real fact, observation, or testimony that would make the conclusion stronger without changing the meaning of the key term.

\textbf{Notice chosen:} \rule{1.8in}{0.4pt}

\textbf{Meaning of the key term you will keep:}\studentline
\textbf{New fact, observation, or testimony:}\studentline\studentline
\textbf{Stronger rewritten argument:}\studentline\studentline\studentline
\textbf{Why does this version give the conclusion better support?}\studentline\studentline

\Needspace{0.42\textheight}
\section*{Part III: Review from Week 1}
Use the Week 1 four-step test on this new claim:

\focusbox{``The Head Tutor wrote the rule in gold ink, so the rule must be wise.''}

\textbf{Conclusion:}\shortline
\textbf{Stated reason:}\shortline
\textbf{Hidden assumption:}\shortline
\textbf{Does the conclusion follow? Explain briefly.}\studentline\studentline

\Needspace{0.38\textheight}
\section*{Part IV: What Can Logic Decide?}
Mark each item as \textbf{Logic can test this} or \textbf{Logic alone cannot decide this}. Then explain two choices.

\begin{enumerate}[leftmargin=*]
\item Whether the word \textit{fair} is being used in two senses. \hfill \rule{2.25in}{0.4pt}
\item Whether the old Head Tutor was morally wise. \hfill \rule{2.25in}{0.4pt}
\item Whether the conclusion follows after \textit{free} is fixed to one meaning. \hfill \rule{2.25in}{0.4pt}
\item Whether Nora actually has a legal right to the cabinet. \hfill \rule{2.25in}{0.4pt}
\item Whether a hidden assumption is needed for the argument to work. \hfill \rule{2.25in}{0.4pt}
\end{enumerate}

\textbf{Explain two choices:}\studentline\studentline\studentline

\Needspace{0.46\textheight}
\section*{Part V: Macbeth Bridge}
Macbeth hears the witches say that ``none of woman born'' shall harm him. He treats the phrase as if it means, ``No human being can harm me.''

\begin{enumerate}[leftmargin=*]
\Needspace{0.12\textheight}\item What meaning does Macbeth assume?\studentline
\Needspace{0.12\textheight}\item What other meaning does the phrase allow?\studentline
\Needspace{0.14\textheight}\item What hidden assumption does Macbeth make about the phrase?\studentline\studentline
\Needspace{0.18\textheight}\item Why is this a dangerous equivocation rather than a simple falsehood?\studentline\studentline
\end{enumerate}

\Needspace{0.18\textheight}
\section*{Final Whately Claim}
In one sentence, explain why a thinker must fix the meaning of key terms before deciding whether an argument follows.

\studentline\studentline

\end{document}
